\documentclass[12pt,a4paper]{article}
\usepackage[utf8]{inputenc}
\usepackage[T1]{fontenc}
\usepackage[romanian]{babel}
\usepackage[margin=2.5cm]{geometry}
\usepackage{graphicx}
\usepackage{listings}
\usepackage{xcolor}
\usepackage{caption}
\usepackage{booktabs}

\newcommand{\studentgroup}{FAF-242}
\definecolor{termbg}{rgb}{0.96,0.96,0.96}
\definecolor{codegray}{rgb}{0.5,0.5,0.5}

\lstset{
  backgroundcolor=\color{termbg},
  basicstyle=\ttfamily\footnotesize,
  breaklines=true,
  captionpos=b,
  frame=single,
  keepspaces=true,
  basewidth={0.52em,0.45em},
  numbers=left,
  numbersep=6pt,
  numberstyle=\tiny\color{codegray},
  showstringspaces=false,
  tabsize=4
}
\lstdefinestyle{term}{numbers=none}

\begin{document}

\begin{titlepage}
    \centering
    \IfFileExists{utm_logo(1).png}{%
        \includegraphics[width=0.25\textwidth]{utm_logo(1).png}\\[0.6cm]%
    }{%
        \vspace*{2cm}%
    }
    {\large \textbf{UNIVERSITATEA TEHNICĂ A MOLDOVEI}}\\[0.2cm]
    {\large Facultatea Calculatoare, Informatică și Microelectronică}\\[0.2cm]
    {\large Departamentul Ingineria Software și Automatică}\\[2.5cm]

    {\Large \textbf{LUCRAREA DE LABORATOR NR. 11, 12 \& 13}}\\[0.4cm]
    {\LARGE \textbf{The Race and the Semaphore \& Producer, Consumer (Linux) \&\\ Semaphores in Your Own Kernel (xv6)}}\\[0.4cm]
    {\large Disciplina: Sisteme de Operare}\\[2.5cm]

    \begin{minipage}{0.45\textwidth}
        \flushleft
        \textbf{Done by:}\\
        st. gr. \studentgroup\\
        Daniel-Ioan \textsc{Strunga}
    \end{minipage}
    \hfill
    \begin{minipage}{0.50\textwidth}
        \flushright
        \textbf{Checked by:}\\
        asist. univ. Patricia \textsc{Reițman}\\
        asist. univ. Marius \textsc{Dumitrașcu}
    \end{minipage}

    \vfill
    {\large Chișinău --- 2026}
\end{titlepage}

\newpage
\tableofcontents
\newpage

\section{Scopul lucrării}
Cele trei laboratoare formează o singură temă: sincronizarea. În Lab 11 (Track A, Linux) se observă o condiție de cursă între două thread-uri care incrementează același contor, apoi se repară cu un mutex și cu un semafor. Același semafor e folosit în trei feluri: ca lock (pornește de la 1), ca semnal (pornește de la 0) și ca numărător (pornește de la N). În Lab 12 (Track A, Linux) instrumentele lucrează împreună în cele două probleme clasice, producător--consumator și cititori--scriitori. Se introduce variabila de condiție, iar la final se construiește și se elimină un deadlock. În Lab 13 (Track B, xv6) semaforul este construit în interiorul kernelului ca trei apeluri de sistem noi (\texttt{sem\_init}, \texttt{sem\_down}, \texttt{sem\_up}, numerele 24--26), dintr-un număr, un spinlock și mecanismul \texttt{sleep}/\texttt{wakeup}.

Toate programele au fost rulate pe laptopul de laborator (Ubuntu 22.04, \texttt{user@user-VPCEH1L0E}). Pentru xv6 a fost nevoie de QEMU 8.2.7, compilat cu scriptul de setup al cursului, deoarece QEMU 6.2 din Ubuntu nu pornește kernelul xv6 actual.

% =========================================================================
\section{Laboratorul 11: The Race and the Semaphore (Linux)}

\subsection{Partea 1 -- cursa (\texttt{race.c})}
\begin{lstlisting}[caption={race.c -- două thread-uri, un contor neprotejat},language=C]
/* race.c - two threads add to the same counter, with no protection.
 * Build: gcc -Wall -pthread -o race race.c       (no -O2: keep counter++ in memory)
 * Run:   ./race   (several times)
 */
#include <stdio.h>
#include <pthread.h>

#define N 1000000

long counter = 0;                          /* shared by both threads */

void *add(void *arg) {
    for (int i = 0; i < N; i++)
        counter++;                         /* read, add 1, write back: three steps */
    return NULL;
}

int main(void) {
    pthread_t a, b;
    pthread_create(&a, NULL, add, NULL);
    pthread_create(&b, NULL, add, NULL);
    pthread_join(a, NULL);
    pthread_join(b, NULL);
    printf("counter = %ld (expected %d)\n", counter, 2 * N);
    return 0;
}
\end{lstlisting}
\begin{lstlisting}[caption={Rezultatul race: niciodată 2\,000\,000}]
===== PART 1: race =====
$ ./race
counter = 1044369 (expected 2000000)
$ ./race
counter = 1093544 (expected 2000000)
$ ./race
counter = 1041945 (expected 2000000)
\end{lstlisting}

\textbf{Observe:}
\begin{enumerate}
    \item \texttt{counter++} înseamnă trei pași: \textbf{load} (citește \texttt{counter} din memorie într-un registru), \textbf{add} (registru + 1) și \textbf{store} (scrie registrul înapoi).
    \item Fie \texttt{counter = 5}. A face load (5), B face load (5), A face add și store $\rightarrow$ 6, B face add și store $\rightarrow$ 6. Au avut loc două incrementări, dar valoarea a crescut cu 1, deci o adunare s-a pierdut.
    \item Felul în care se intercalează thread-urile depinde de planificator, de nucleele CPU și de timing, care diferă la fiecare rulare (1\,044\,369, 1\,093\,544, 1\,041\,945).
\end{enumerate}

\subsection{Partea 2 -- mutex, apoi semafor}
Singura diferență față de \texttt{race.c} este protecția din jurul lui \texttt{counter++}:
\begin{lstlisting}[language=C,caption={mutex.c -- sectiunea critica}]
pthread_mutex_lock(&lock);      /* only one thread past this line */
counter++;                      /* the critical section */
pthread_mutex_unlock(&lock);    /* let the other one in */
\end{lstlisting}
\begin{lstlisting}[language=C,caption={sem.c -- acelasi lock, din semafor (sem\_init(\&s, 0, 1))}]
sem_wait(&s);                   /* 1 -> 0: I go in; at 0: wait */
counter++;
sem_post(&s);                   /* 0 -> 1: the next one may go in */
\end{lstlisting}
\begin{lstlisting}[caption={mutex, sem, sem cu start 2 și timpii de execuție}]
===== PART 2: mutex si sem =====
$ ./mutex
counter = 2000000 (expected 2000000)
$ ./mutex
counter = 2000000 (expected 2000000)
$ ./mutex
counter = 2000000 (expected 2000000)
$ ./sem
counter = 2000000 (expected 2000000)
$ ./sem
counter = 2000000 (expected 2000000)
$ ./sem
counter = 2000000 (expected 2000000)
--- sem_init(&s, 0, 2) ---
$ ./sem
counter = 1990835 (expected 2000000)
$ ./sem
counter = 1937654 (expected 2000000)
$ ./sem
counter = 1921909 (expected 2000000)
--- time ---
$ time ./race
counter = 1031595 (expected 2000000)

real	0m0.014s
user	0m0.025s
sys	0m0.001s
$ time ./mutex
counter = 2000000 (expected 2000000)

real	0m0.196s
user	0m0.390s
sys	0m0.001s
$ time ./sem
counter = 2000000 (expected 2000000)

real	0m0.402s
user	0m0.468s
sys	0m0.331s
\end{lstlisting}

\textbf{Observe:}
\begin{enumerate}
    \item Da: \texttt{mutex} și \texttt{sem} dau 2\,000\,000 la fiecare rulare (3 din 3).
    \item Cu \texttt{sem\_init(\&s, 0, 2)}, semaforul lasă \textbf{două} thread-uri deodată în secțiunea critică, deci cursa revine (1\,990\,835, 1\,937\,654, 1\,921\,909). Pierderile sunt mai mici decât la \texttt{race} doar pentru că \texttt{sem\_wait}/\texttt{sem\_post} încetinesc bucla.
    \item Fiecare incrementare costă în plus un lock și un unlock (operații atomice). Când thread-ul celălalt ține lock-ul, se intră în kernel (\texttt{futex}) și se așteaptă, așa că thread-urile ajung să lucreze pe rând. Pe vaio: race 0.014\,s (greșit), mutex 0.196\,s, sem 0.402\,s, din care 0.331\,s petrecute în kernel.
\end{enumerate}

\subsection{Partea 3 -- semaforul ca semnal (\texttt{order.c})}
\begin{lstlisting}[language=C,caption={order.c -- B asteapta semnalul lui A}]
sem_t ready;                 /* sem_init(&ready, 0, 0): inca nimic gata */
void *writer(void *arg) { sleep(1); strcpy(message, "hello");
                          printf("A: message written\n"); sem_post(&ready); return NULL; }
void *reader(void *arg) { printf("B: waiting for the message...\n");
                          sem_wait(&ready); printf("B: I read \"%s\"\n", message); return NULL; }
\end{lstlisting}
\begin{lstlisting}[caption={order cu start 0 și cu start 1}]
===== PART 3: order =====
$ ./order
B: waiting for the message...
A: message written
B: I read "hello"
--- sem_init(&ready, 0, 1) ---
$ ./order
B: waiting for the message...
B: I read ""
A: message written
\end{lstlisting}

\textbf{Observe:}
\begin{enumerate}
    \item Cu start 1, B citește \texttt{""} (șir gol): \texttt{sem\_wait} găsește valoarea 1 și trece imediat, înainte ca A să fi scris mesajul.
    \item Un semafor care pornește de la 0 înseamnă „încă nu s-a întâmplat nimic; așteaptă până dă cineva semnalul cu \texttt{sem\_post}”.
\end{enumerate}

\subsection{Partea 4 -- semaforul ca numărător (\texttt{parking.c})}
\begin{lstlisting}[caption={parking cu 3 locuri și cu 1 loc}]
===== PART 4: parking =====
$ ./parking
car 1 enters (inside: 1)
car 2 enters (inside: 2)
car 3 enters (inside: 3)
car 1 leaves (inside: 2)
car 2 leaves (inside: 1)
car 4 enters (inside: 2)
car 5 enters (inside: 3)
car 3 leaves (inside: 2)
car 6 enters (inside: 3)
car 4 leaves (inside: 2)
car 6 leaves (inside: 1)
car 5 leaves (inside: 0)
--- PLACES 1 ---
$ ./parking
car 1 enters (inside: 1)
car 1 leaves (inside: 0)
car 2 enters (inside: 1)
car 2 leaves (inside: 0)
car 3 enters (inside: 1)
car 3 leaves (inside: 0)
car 4 enters (inside: 1)
car 4 leaves (inside: 0)
car 5 enters (inside: 1)
car 5 leaves (inside: 0)
car 6 enters (inside: 1)
car 6 leaves (inside: 0)
\end{lstlisting}

\textbf{Observe:}
\begin{enumerate}
    \item Valoarea maximă a lui \texttt{inside} este \textbf{3}. O garantează \texttt{sem\_wait(\&places);}, cu \texttt{sem\_init(\&places, 0, PLACES)} și \texttt{PLACES = 3}: a patra mașină doarme până pleacă una.
    \item Cu \texttt{PLACES = 1} semaforul devine binar, adică un \textbf{mutex}: mașinile intră strict pe rând.
\end{enumerate}
\textbf{Opțional (atomic):} \texttt{atomic\_fetch\_add} dă 2\,000\,000 în 0.079\,s, fără niciun lock, pentru că load, add și store se fac ca o singură instrucțiune indivizibilă.

\newpage
% =========================================================================
\section{Laboratorul 12: Producer, Consumer and the Other Tools (Linux)}

\subsection{Partea 1 -- producător și consumator}
\begin{lstlisting}[caption={prodcons.c -- producătorul (empty = sloturi libere, full = sloturi ocupate)},language=C]
void *producer(void *arg) {
    for (int item = 1; item <= ITEMS; item++) {
        sem_wait(&empty);               /* wait for a free slot */
        pthread_mutex_lock(&lock);
        buf[in] = item;
        in = (in + 1) % SIZE;
        count++;
        printf("put %d    (in buffer: %d)\n", item, count);
        pthread_mutex_unlock(&lock);
        sem_post(&full);                /* one more filled slot */
    }
    return NULL;
}
\end{lstlisting}
\begin{lstlisting}[caption={prodcons normal și cu lock înainte de sem\_wait(\&empty)}]
===== PART 1: prodcons =====
$ ./prodcons
put 1    (in buffer: 1)
put 2    (in buffer: 2)
put 3    (in buffer: 3)
put 4    (in buffer: 4)
            take 1   (in buffer: 3)
put 5    (in buffer: 4)
            take 2   (in buffer: 3)
put 6    (in buffer: 4)
            take 3   (in buffer: 3)
put 7    (in buffer: 4)
            take 4   (in buffer: 3)
put 8    (in buffer: 4)
            take 5   (in buffer: 3)
            take 6   (in buffer: 2)
            take 7   (in buffer: 1)
            take 8   (in buffer: 0)
--- lock inainte de sem_wait(&empty) ---
$ timeout 5 ./prodcons
put 1    (in buffer: 1)
put 2    (in buffer: 2)
put 3    (in buffer: 3)
put 4    (in buffer: 4)
exit code: 124
\end{lstlisting}

\textbf{Observe:}
\begin{enumerate}
    \item Producătorul este oprit de \texttt{sem\_wait(\&empty)}: după 4 elemente \texttt{empty} ajunge la 0.
    \item La început așteaptă \textbf{consumatorul}, în \texttt{sem\_wait(\&full)}, pentru că \texttt{full} pornește de la 0.
    \item Cu lock-ul luat înaintea lui \texttt{sem\_wait(\&empty)}, programul se blochează după \texttt{put 4} (exit 124). Producătorul ține mutex-ul și doarme pe \texttt{empty}, iar consumatorul nu poate lua mutex-ul ca să elibereze un slot. Fiecare îl așteaptă pe celălalt: \textbf{deadlock}.
\end{enumerate}

\subsection{Partea 2 -- variabila de condiție}
\begin{lstlisting}[caption={condvar.c -- cititorul},language=C]
void *reader(void *arg) {
    pthread_mutex_lock(&lock);
    while (!ready) {                            /* while, not if: re-check after waking */
        printf("B: not ready, I sleep\n");
        pthread_cond_wait(&changed, &lock);     /* unlocks, sleeps, locks again */
    }
    printf("B: I read \"%s\"\n", message);
    pthread_mutex_unlock(&lock);
    return NULL;
}
\end{lstlisting}
\begin{lstlisting}[caption={condvar cu while și cu if}]
===== PART 2: condvar =====
$ ./condvar
B: not ready, I sleep
A: message written
B: I read "hello"
--- if (!ready) ---
$ ./condvar
B: not ready, I sleep
A: message written
B: I read "hello"
\end{lstlisting}

\textbf{Observe:}
\begin{enumerate}
    \item Starea „mesaj gata” e ținută minte de variabila \texttt{ready}, protejată de mutex. Variabila de condiție nu ține minte nimic: un \texttt{signal} fără nimeni care să aștepte se pierde.
    \item Cu \texttt{if (!ready)} merge și aici (output identic): există un singur cititor și un singur semnal.
    \item \texttt{pthread\_cond\_wait} se poate întoarce și fără semnal (\emph{spurious wakeup}). În plus, cu mai mulți cititori sau cu \texttt{broadcast}, alt thread poate lua mutex-ul primul și poate schimba condiția. \texttt{while} re-verifică \texttt{ready} după fiecare trezire.
\end{enumerate}

\subsection{Partea 3 -- cititori și scriitori}
\begin{lstlisting}[caption={readers.c -- cititorul},language=C]
void *reader(void *arg) {
    long id = (long)arg;
    sem_wait(&rlock);
    if (++readers == 1)         /* the first reader in locks out the writers */
        sem_wait(&room);
    sem_post(&rlock);

    printf("reader %ld reads %d   (readers inside: %d)\n", id, value, readers);
    usleep(200000);

    sem_wait(&rlock);
    if (--readers == 0)         /* the last reader out lets the writers in */
        sem_post(&room);
    sem_post(&rlock);
    return NULL;
}
\end{lstlisting}
\begin{lstlisting}[caption={readers: trei cititori împreună, scriitorii singuri}]
===== PART 3: readers =====
$ ./readers
writer 1 writes 10 (alone)
reader 1 reads 10   (readers inside: 1)
reader 2 reads 10   (readers inside: 2)
reader 3 reads 10   (readers inside: 3)
writer 2 writes 20 (alone)
\end{lstlisting}

\textbf{Observe:}
\begin{enumerate}
    \item Liniile \texttt{if (++readers == 1) sem\_wait(\&room);} (primul cititor închide camera) și \texttt{if (--readers == 0) sem\_post(\&room);} (ultimul o deschide). Cititorii din mijloc nu ating \texttt{room}.
    \item \texttt{++readers}/\texttt{--readers} sunt read-modify-write, adică aceeași cursă ca în \texttt{race.c}. Fără \texttt{rlock}, doi cititori ar putea vedea amândoi 1 (dublu \texttt{sem\_wait(\&room)}, deci blocaj) sau amândoi 0 (camera eliberată cât încă citește cineva, iar scriitorul intră peste cititori).
    \item Dacă vin mereu cititori noi, \texttt{readers} nu ajunge niciodată la 0 și writer 2 așteaptă la nesfârșit: \textbf{starvation}.
\end{enumerate}

\subsection{Partea 4 -- deadlock, apoi fix}
\begin{lstlisting}[caption={deadlock: original, fără sleep (de 10 ori), apoi după fix}]
===== PART 4: deadlock =====
$ timeout 5 ./deadlock; echo "exit code: $?"
one: I have A, I want B
two: I have B, I want A
exit code: 124
--- fara sleep(1), de 10 ori ---
run 1:
one: I have A, I want B
two: I have B, I want A
exit code: 124
run 2:
one: I have A, I want B
two: I have B, I want A
exit code: 124
run 3:
one: I have A, I want B
two: I have B, I want A
exit code: 124
run 4:
one: I have A, I want B
two: I have B, I want A
exit code: 124
run 5:
one: I have A, I want B
two: I have B, I want A
exit code: 124
run 6:
one: I have A, I want B
two: I have B, I want A
exit code: 124
run 7:
one: I have A, I want B
two: I have B, I want A
exit code: 124
run 8:
one: I have A, I want B
two: I have B, I want A
exit code: 124
run 9:
one: I have A, I want B
two: I have B, I want A
exit code: 124
run 10:
one: I have A, I want B
two: I have B, I want A
exit code: 124
--- dupa fix (A apoi B si in two) ---
$ timeout 5 ./deadlock; echo "exit code: $?"
one: I have A, I want B
one: I have both
two: I have B, I want A
two: I have both
done
exit code: 0
\end{lstlisting}
\begin{lstlisting}[caption={deadlock.c -- thread-ul two după fix (A, apoi B, ca în one)},language=C]
void *two(void *arg) {
    pthread_mutex_lock(&A);             /* the fix: A first, same order as one() */
    printf("two: I have B, I want A\n");
    sleep(1);
    pthread_mutex_lock(&B);             /* then B */
    printf("two: I have both\n");
    pthread_mutex_unlock(&B);
    pthread_mutex_unlock(&A);
    return NULL;
}
\end{lstlisting}

\textbf{Observe:}
\begin{enumerate}
    \item \texttt{one} ține \textbf{A} și așteaptă \textbf{B}, iar \texttt{two} ține \textbf{B} și așteaptă \textbf{A}. Cercul: one $\rightarrow$ B (ținut de two) $\rightarrow$ A (ținut de one).
    \item Pe vaio, fără \texttt{sleep}, s-a blocat de 10 ori din 10, pentru că \texttt{printf} dintre cele două lock-uri lasă suficient timp. În general rezultatul depinde de timing; pe altă mașină s-a blocat doar de 3 ori din 10. E mai rău, nu mai bine: bug-ul devine intermitent, trece de teste și apare rar, greu de reprodus.
    \item Dacă toate thread-urile iau lock-urile în aceeași ordine, niciunul nu poate ține B așteptând A, deci cercul de așteptare nu se poate forma.
\end{enumerate}

\newpage
% =========================================================================
\section{Laboratorul 13: Semaphores in Your Own Kernel (xv6)}

\subsection{Partea 1 -- numerele și partea user}
Au fost adăugate \texttt{SYS\_sem\_init 24}, \texttt{SYS\_sem\_down 25} și \texttt{SYS\_sem\_up 26} în \texttt{kernel/syscall.h}, declarațiile în \texttt{user/user.h}, stub-urile în \texttt{user/usys.pl} și programele \texttt{semtest} și \texttt{mix} în \texttt{UPROGS}. Fără partea de kernel, apelurile întorc $-1$ și nu blochează nimic, deci copilul rulează înaintea părintelui:
\begin{lstlisting}[caption={Check 1 -- unknown sys call 24/25/26, child înainte de parent}]
# Check 1: doar partea user (sem_init/sem_down/sem_up necunoscute in kernel)
$ semtest
3 semtest: unknown sys call 24
4 semtest: unknown sys call 25
child: I ran after the parent
parent: done, I let the child go
3 semtest: unknown sys call 26
3 semtest: unknown sys call 24
3 semtest: unknown sys call 24
3 semtest: unknown sys call 25
p5 semtest: unknown sys call 25
i ng  p1o
n3 semtest: unknown sys call 26
g3 semtest: unknown sys call 25
pin g 12

5 semtest: unknown sys call 26
3 semtest: unknown sys call 26
5 semtest: unknown sys call 25
 3 semtest: unknown sys call 25
 popnign 2g 3

3 semtest: unknown sys call 26
5 semtest: unknown sys call 26
5 semtest: unknown sys call 25
  pong 3
5 semtest: unknown sys call 26
semtest: OK
\end{lstlisting}


\subsection{Partea 2 -- partea de kernel (\texttt{kernel/sem.c})}
Un semafor este un număr plus un spinlock care îl protejează. Tabela are 8 semafoare, comune tuturor proceselor. \texttt{sem.o} a fost adăugat în \texttt{OBJS}, iar cele trei funcții în tabela \texttt{syscalls[]} din \texttt{kernel/syscall.c}.
\begin{lstlisting}[caption={kernel/sem.c -- sem\_down},language=C]
// sem_down(id): take one unit; if there is none, sleep until there is.
uint64
sys_sem_down(void)
{
  int id;
  argint(0, &id);
  if (id < 0 || id >= NSEM)
    return -1;
  struct sem *s = &sems[id];

  acquire(&s->lock);
  while (s->value == 0) {
    sleep_prepare(s);       // I will sleep on the address of s
    release(&s->lock);      // never sleep holding a spinlock
    sleep();                // returns at once if sem_up came in between
    acquire(&s->lock);
    if (killed(myproc())) { // kill must still work on a waiting process
      release(&s->lock);
      return -1;
    }
  }
  s->value--;
  release(&s->lock);
  return 0;
}
\end{lstlisting}
\begin{lstlisting}[caption={kernel/sem.c -- sem\_up},language=C]
// sem_up(id): give one unit back and wake whoever sleeps on s.
uint64
sys_sem_up(void)
{
  int id;
  argint(0, &id);
  if (id < 0 || id >= NSEM)
    return -1;
  struct sem *s = &sems[id];

  acquire(&s->lock);
  s->value++;
  wakeup(s);
  release(&s->lock);
  return 0;
}
\end{lstlisting}
\begin{lstlisting}[caption={Check 2 -- semtest: OK}]
# Check 2: cu kernel/sem.c
$ semtest
parent: done, I let the child go
child: I ran after the parent
ping 1
  pong 1
ping 2
  pong 2
ping 3
  pong 3
semtest: OK
\end{lstlisting}

Părintele așteaptă 10 tick-uri înainte să-l elibereze pe copil, iar copilul afișează al doilea pentru că dormea în \texttt{sem\_down}. Ping și pong alternează strict.

\subsection{Partea 3 -- semaforul ca lock: \texttt{mix}}
\begin{lstlisting}[caption={mix fără lock -- literele se amestecă}]
# Part 3: mix fara lock (de doua ori)
$ mix
AbAbbAbAbAbAbAbAbbbbAbbbbbbAbAbAbAbAbAbAbAbAAbbbbA
AbAbAbbbAbAbAbAbAbAbbbAb
bbAbAbAbAbAAAbAbAbAAbAbAbAbbbbA
AAAAbAbAAbAbAbAbAbAbAbAbbA
AAAbAbAbAbAbAbAbAAbAbAbAbAbbAbAbAbbbAbA
AbAbAbAbAbAbAAAAb
bAbAbAbAbAbbAbAAAbAbAbAbAbAbAAAbAAAAAbAbAbAbbAbAAbAbA

bAbAbAbAbAbAbAbAbAbAbAbAbAAbAbAAAAbbAbAbAAAbAAbAbAbAbAb
bbbbb
$ mix
AbAbAbAbAbbbAbAbAbAbAbAbAbAbAbAbAbAbAbAbAbAbAbAbAAAbAbAbAbAb
bb
bAbAbAbAbAbAbAbAbAbAbAbAbAbAbbAbAbAbAbAbAbAbAbAbAbAbAbA
AbAbAb
bbAbbAbAbAbAbAbAbAbAbAbAbAbAbbbbAAAAbAbAbAbAbAbAbAb
bAbAbAbAbAbAb
bAbbAbbbbAbAbAbAbAbbbAbAbAAbAbAbAbAbAbAbA
AbAbAbAbAbAbAbAbAbAbAbAb
bAbAbAbAbAbAbAbbAbAbAbAbAAAbAbAbAbAbA
AAAAAAAAAAA
\end{lstlisting}
\begin{lstlisting}[caption={mix lock -- fiecare linie rămâne întreagă}]
# Part 3: mix cu lock
$ mix lock
AAAAAAAAAAAAAAAAAAAAAAAAAAAAAA
AAAAAAAAAAAAAAAAAAAAAAAAAAAAAA
AAAAAAAAAAAAAAAAAAAAAAAAAAAAAA
AAAAAAAAAAAAAAAAAAAAAAAAAAAAAA
AAAAAAAAAAAAAAAAAAAAAAAAAAAAAA
bbbbbbbbbbbbbbbbbbbbbbbbbbbbbb
bbbbbbbbbbbbbbbbbbbbbbbbbbbbbb
bbbbbbbbbbbbbbbbbbbbbbbbbbbbbb
bbbbbbbbbbbbbbbbbbbbbbbbbbbbbb
bbbbbbbbbbbbbbbbbbbbbbbbbbbbbb
\end{lstlisting}

\textbf{Observe (notes.txt):}
\begin{enumerate}
    \item \texttt{mix} scrie câte o literă per \texttt{write()}, deci o linie costă 31 de apeluri de sistem. Între oricare două apeluri, celălalt proces (pe alt CPU, din 3) poate scrie și el, așa că literele se intercalează.
    \item \texttt{sem\_down(3);} înainte de bucla care scrie linia și \texttt{sem\_up(3);} după ea. Semaforul 3 pornește de la 1, deci funcționează ca un lock.
    \item \texttt{sem\_up} face \texttt{wakeup(s)}, care trezește \emph{toate} procesele care dorm pe adresa \texttt{s}, dar adaugă o singură unitate. Primul care ia spinlock-ul o consumă, iar ceilalți găsesc din nou 0 și trebuie să se culce la loc.
    \item \texttt{sem\_up} trebuie să ia \texttt{s->lock} ca să incrementeze și să trezească. Dacă procesul ar dormi ținând spinlock-ul, \texttt{sem\_up} s-ar învârti la nesfârșit în \texttt{acquire}: deadlock. În plus, un spinlock ținut înseamnă întreruperi dezactivate, iar \texttt{sched()} din xv6 dă \texttt{panic("sched locks")}. \texttt{sleep\_prepare(s)} înainte de \texttt{release} face ca un \texttt{sem\_up} venit între \texttt{release} și \texttt{sleep} să nu se piardă.
\end{enumerate}

\newpage
% =========================================================================
\section{Concluzii}
\begin{enumerate}
    \item \textbf{Cursa e reală și se vede:} o operație care arată ca un singur pas (\texttt{counter++}) are de fapt trei pași. Pe vaio, două thread-uri au pierdut aproape jumătate din cele 2\,000\,000 de adunări, iar rezultatul a diferit la fiecare rulare.
    \item \textbf{Un singur instrument, trei roluri:} semaforul pornit de la 1 este un lock, cel pornit de la 0 este un semnal, iar cel pornit de la N numără resurse. Valoarea de start decide comportamentul (start 2 strică lock-ul, start 1 strică semnalul).
    \item \textbf{Instrumentele lucrează împreună și pot greși împreună:} în producător--consumator și în cititori--scriitori, ordinea apelurilor contează. Lock-ul luat înaintea lui \texttt{sem\_wait}, sau două mutex-uri luate în ordini diferite, duc la deadlock. Regula care îl previne este ordinea globală a lock-urilor. Variabila de condiție cere starea într-o variabilă separată și \texttt{while}, nu \texttt{if}.
    \item \textbf{În kernel, semaforul e mic:} un număr, un spinlock și \texttt{sleep}/\texttt{wakeup} pe adresa semaforului. Procesele care așteaptă dorm cu adevărat, fără să consume CPU. Detaliile delicate sunt bucla \texttt{while} și eliberarea spinlock-ului înainte de \texttt{sleep()}, exact ca la variabila de condiție din Linux.
\end{enumerate}

\end{document}
